\section{נספח: מילון מונחים \en{(Glossary)}}

\begin{table}[h]
\centering
\caption{מונחים מרכזיים במסמך}
\label{tab:glossary}
\small
\begin{tabularx}{\textwidth}{>{\bfseries\en}l X}
\toprule
\theadrow \thead{\en{Term}} & \thead{הסבר} \\
\midrule
PSNR & יחס אות לרעש שיא, בדציבלים; מודד כמה התמונה המשוחזרת קרובה למקור \\
SSIM & מדד דמיון מבני בין שתי תמונות (בהירות, ניגודיות, מבנה) \\
GPU & מעבד גרפי המאיץ הרצה של מודלי למידה עמוקה \\
Super-resolution & הגדלת רזולוציה של תמונה תוך שחזור פרטים \\
Break-even & נקודת איזון: היקף המכירות שבו הרווח התפעולי מגיע לאפס \\
Contribution margin & תרומה להזמנה: מחיר פחות עלות משתנה \\
Monte Carlo & שיטה לחישוב תוצאות בהרצה חוזרת עם אי־ודאות אקראית \\
BiDi & דו־כיווניות: ערבוב של כתיבה מימין לשמאל ומשמאל לימין \\
CLS & קובץ תבנית של \en{LaTeX} הקובע את עיצוב המסמך \\
Skill & קובץ הוראות שמגדיר תפקיד וכללי עבודה לסוכן \en{AI} \\
\bottomrule
\end{tabularx}
\end{table}
